
% BEGIN S2
 
\par\medskip\noindent\textbf{2 Ultrarings}\par
 
% BEGIN S2.SS1
 
\par\medskip\noindent\textbf{2.1 Main Definitions}\par
 
% BEGIN S2.SS1.p1
 
In what follows, we often abbreviate composition of arrows $f\circ g$ by $fg$. If $(\mathcal{C},\otimes,I,\gamma)$ is a symmetric monoidal category, where $\gamma$ is its symmetry, we denote by $\beta:I\otimes I\to I$ the canonical isomorphism. By “symmetric monoidal functor” we will always mean strong symmetric monoidal functor (i.e., tensor and unit are preserved up to iso).\par
 
% END S2.SS1.p1
 
% BEGIN S2.SS1.p2
 
A Frobenius sesquimonoid in a symmetric monoidal category is an object $A$ equipped with arrows $\delta:A\to A\otimes A$, $\varepsilon:A\to I$, $\delta^{\ast}:A\otimes A\to A$ such that $(A,\delta,\varepsilon)$ is a commutative comonoid, $\delta^{\ast}$ is commutative ($\delta^{\ast}\gamma_{A,A}=\delta^{\ast}$), associative ($\delta^{\ast}(\mathrm{id}\otimes\delta^{\ast})=\delta^{\ast}(\delta^{\ast}%
\otimes\mathrm{id})$, left inverse to $\delta$ ($\delta^{\ast}\delta=\mathrm{id}_{A}$) and verifies the Frobenius law $(\mathrm{id}_{A}\otimes\delta^{\ast})(\delta\otimes\mathrm{id}_{A})=\delta%
\delta^{\ast}=(\delta^{\ast}\otimes\mathrm{id}_{A})(\mathrm{id}_{A}\otimes\delta)$.\par
 
% END S2.SS1.p2
 
% BEGIN S2.SS1.p3
 
Initial objects will be denoted by $0$. Recall that a (binary) coproduct diagram over an object $A$ is a pair of arrows $\iota_{1}:A_{1}\to A$, $\iota_{2}:A_{2}\to A$, called injections, such that, for every $f_{1}:A_{1}\to B$, $f_{2}:A_{2}\to B$ there is a unique arrow $[f_{1},f_{2}]:A\to B$ such that $[f_{1},f_{2}]\iota_{i}=f_{i}$. In that case, we write $A=A_{1}\oplus A_{2}$.\par
 
% END S2.SS1.p3
 
% BEGIN S2.SS1.p4
 
A distributive monoidal category is a symmetric monoidal category with finite coproducts such that, for all objects $A,B,C$, $A\otimes 0$ is initial and the arrow $\chi_{A,B,C}:=[\mathrm{id}_{A}\otimes\iota_{1},\mathrm{id}_{A}\otimes\iota_{2}%
]:(A\otimes B)\oplus(A\otimes C)\to A\otimes(B\oplus C)$ is an isomorphism.\par
 
% END S2.SS1.p4
 
% BEGIN S2.SS1.p5
 
Recall that an initial object which is also terminal is called a zero object and that any category with a zero object has, for any pair of objects $A,B$, a zero morphism $0_{A,B}:A\to B$ obtained by composing the terminal and initial arrows $A\to 0\to B$. Also, for any pair of objects $A_{1},A_{2}$, the category has coinjections222We refer to these as coinjections rather than projections to avoid the reader assuming that the coproducts are also products. defined by:\par
    $\displaystyle\iota_{1}^{\ast}$ $\displaystyle:=[\mathrm{id}_{A_{1}},0_{A_{2},A_{1}}]:A_{1}\oplus A_{2}\to A_{1},$ $\displaystyle\iota_{2}^{\ast}$ $\displaystyle:=[0_{A_{1},A_{2}},\mathrm{id}_{A_{2}}]:A_{1}\oplus A_{2}\to A_{2}.$    
It immediately follows that these verify\par
    \[\iota^{\ast}_{i}\iota_{j}=\left\{\begin{array}[]{ll}\mathrm{id}_{A_{i}}&\text{%
if }i=j\\
0_{A_{j},A_{i}}&\text{if }i\neq j.\end{array}\right.\]  (1)   
% END S2.SS1.p5
 
% BEGIN Thmtheorem1
 
\par\medskip\noindent\textbf{Definition 1.}\par
 
% BEGIN Thmtheorem1.p1
 
A $\delta^{\ast}$-category is a distributive monoidal category $(\mathcal{C},\otimes,I,\gamma)$ with a zero object such that each object $A$ is equipped with a Frobenius sesquimonoid structure $\delta_{A},\varepsilon_{A},\delta_{A}^{\ast}$ satisfying, for every objects $A,B,A_{1},A_{2}$,\par
    $\displaystyle\delta_{A\otimes B}$ $\displaystyle=(\mathrm{id}_{A}\otimes\gamma_{A,B}\otimes\mathrm{id}_{B})(%
\delta_{A}\otimes\delta_{B})$ $\displaystyle\varepsilon_{A\otimes B}$ $\displaystyle=\beta(\varepsilon_{A}\otimes\varepsilon_{B})$ $\displaystyle\delta_{I}$ $\displaystyle=\beta^{-1},$     $\displaystyle\delta^{\ast}_{A\otimes B}$ $\displaystyle=(\delta^{\ast}_{A}\otimes\delta^{\ast}_{B})(\mathrm{id}_{A}%
\otimes\gamma_{A,B}\otimes\mathrm{id}_{B})$ $\displaystyle\varepsilon_{I}$ $\displaystyle=\mathrm{id}_{I}$ $\displaystyle\delta^{\ast}_{I}$ $\displaystyle=\beta,$     $\displaystyle\hskip-15.0pt\delta^{\ast}_{A_{1}\oplus A_{2}}(\iota_{i}\otimes%
\iota_{j})$ $\displaystyle=\left\{\begin{array}[]{ll}\iota_{i}\delta^{\ast}_{A_{i}}&\text{%
if }i=j\\
0_{A_{i}\otimes A_{j},A_{1}\oplus A_{2}}&\text{if }i\neq j\end{array}\right.$ $\displaystyle\delta_{A_{1}\oplus A_{2}}\iota_{i}$ $\displaystyle=(\iota_{i}\otimes\iota_{i})\delta_{A_{i}}$ $\displaystyle\varepsilon_{A_{1}\oplus A_{2}}\iota_{i}$ $\displaystyle=\varepsilon_{A_{i}},$    
and such that, for all $f,g:A\rightrightarrows B$, $\delta^{\ast}_{B}(f\otimes\mathrm{id})=\delta^{\ast}_{B}(g\otimes\mathrm{id})$ implies $f=g$.\par
 
% END Thmtheorem1.p1
 
% BEGIN Thmtheorem1.p2
 
A morphism $F:\mathcal{C}\to\mathcal{D}$ of $\delta^{\ast}$-categories is a finite-coproduct-preserving symmetric monoidal functor with structural isomorphisms $\varphi_{A,B}:FA\otimes FB\to F(A\otimes B)$ and $\psi:I\to FI$ such that, for each object $A$ of $\mathcal{C}$, $F\delta_{A}=\varphi_{A,A}\delta_{FA}$, $F\delta^{\ast}_{A}=\delta^{\ast}_{FA}\varphi^{-1}_{A,A}$ and $F\varepsilon_{A}=\psi\varepsilon_{FA}$.\par
 
% END Thmtheorem1.p2
 
% END Thmtheorem1
 
% BEGIN S2.SS1.p6
 
In the sequel, subscripts such as those in $\mathrm{id}_{A},\delta_{A},\varepsilon_{A}\ldots$ will be dropped when unambiguously retrievable from the context.\par
 
% END S2.SS1.p6
 
% BEGIN Thmtheorem2
 
\par\medskip\noindent\textbf{Lemma 2.}\par
 
% BEGIN Thmtheorem2.p1
 
In a $\delta^{\ast}$-category, coinjections are jointly monic.\par
 
% END Thmtheorem2.p1
 
% END Thmtheorem2
 
% BEGIN S2.SS1.SSS0.P0.SPx1
 
\par\medskip\noindent\textbf{Proof.}\par
 
% BEGIN S2.SS1.SSS0.P0.SPx1.p1
 
See §A.1. $\hfill\blacktriangleleft$\par
 
% END S2.SS1.SSS0.P0.SPx1.p1
 
% BEGIN S2.SS1.SSS0.P0.SPx1.p2
 
By the above lemma, for every arrows $f_{1}:B\to A_{1}$, $f_{2}:B\to A_{2}$ of a $\delta^{\ast}$-category, there is at most one arrow $\langle f_{1},f_{2}\rangle:B\to A_{1}\oplus A_{2}$, which we call pairing, such that $\iota^{\ast}_{i}\langle f_{1},f_{2}\rangle=f_{i}$. Some pairings automatically exist. For instance, by (1), $\iota_{1}:A_{1}\to A_{1}\oplus A_{2}$ is the pairing $\langle\mathrm{id}_{A_{1}},0_{A_{1},A_{2}}\rangle$ and dually for $\iota_{2}$. The pairing $\langle 0_{B,A_{1}},0_{B,A_{2}}\rangle$ is simply $0_{B,A_{1}\oplus A_{2}}$. More subtly, the inverse $\chi^{-1}_{A,B_{1},B_{2}}$ of the distributivity isomorphism is the pairing $\langle\mathrm{id}_{A}\otimes\iota^{*}_{1},\mathrm{id}_{A}\otimes\iota^{*}_{2}\rangle$, which can be deduced using the fact that injections into a coproduct are jointly epic.\par
 
% END S2.SS1.SSS0.P0.SPx1.p2
 
% BEGIN Thmtheorem3
 
\par\medskip\noindent\textbf{Definition 3.}\par
 
% BEGIN Thmtheorem3.p1
 
An ultraring is a small $\delta^{\ast}$-category such that:\par
  
\par\noindent  1. 
% BEGIN S2.I1.i1.p1
 
disjoint pairings: for every $f:C\to A$ and $g:C\to B$, if $(f\otimes g)\delta=0$ then their pairing $\langle f,g\rangle$ exists;\par
 
% END S2.I1.i1.p1
 
 
\par\noindent  2. 
% BEGIN S2.I1.i2.p1
 
complements: for every $p:A\to I$ such that $(p\otimes p)\delta=p$, there exist a unique $\overline{p}:A\to I$ such that $(p\otimes\overline{p})\delta=0$ and $\nabla\langle p,\overline{p}\rangle=\varepsilon$, where $\nabla:I\oplus I\to I$ is the codiagonal.\par
 
% END S2.I1.i2.p1
 
  
% END Thmtheorem3.p1
 
% BEGIN Thmtheorem3.p2
 
A morphism of ultrarings is just a morphism of $\delta^{\ast}$-categories. We denote by $\mathbf{URing}$ the strict $2$-category whose objects are ultrarings, whose arrows are morphisms between them and whose $2$-cells are monoidal natural transformations.\par
 
% END Thmtheorem3.p2
 
% END Thmtheorem3
 
% END S2.SS1.SSS0.P0.SPx1
 
% END S2.SS1
 
% BEGIN S2.SS2
 
\par\medskip\noindent\textbf{2.2 Examples of Ultrarings}\par
 
% BEGIN S2.SS2.p1
 
Let $R$ be a commutative ring. We denote by $\mathcal{U}R$ the category of free $R$-modules of finite rank, and linear maps between them. This may be equivalently described as the category whose objects are natural numbers and whose arrows $n\to m$ are $m\times n$ matrices ($m$ rows, $n$ columns) with coefficients in $R$, and where composition is matrix multiplication.\par
 
% END S2.SS2.p1
 
% BEGIN S2.SS2.p2
 
This category is distributive monoidal: the (strict) monoidal structure is given by multiplication of natural numbers for objects and by the standard tensor product of matrices (also known as Kronecker product) for morphisms. Coproduct is addition and $0$ is easily seen to be a zero object. The Frobenius sesquimonoid structure is as follows: $\delta_{n}$ is the $n^{2}\times n$ matrix whose $i$-th column has the element $1$ at position $n(i-1)+i$ and is zero everywhere else; $\varepsilon_{n}$ is the row matrix $(1\ldots 1)$; and $\delta^{\ast}_{n}$ is the transpose of $\delta_{n}$.\par
 
% END S2.SS2.p2
 
% BEGIN S2.SS2.p3
 
Under matrix addition, $\mathcal{U}R$ is actually an additive category, which means it has biproducts, so point (1) of Definition 3 is immediate. For what concerns point (2), any $p:n\to 1$ is a row matrix $(p_{1}\ldots p_{n})$, and the condition $(p\otimes p)\delta=p$ amounts to $p_{i}^{2}=p_{i}$ for all $i$. It is then a matter of elementary calculations to show that $\overline{p}:=(1-p_{1}\ldots 1-p_{n})$ is the unique morphism $n\to 1$ verifying $(p\otimes\overline{p})\delta=0$ and $\nabla\langle p,\overline{p}\rangle=\varepsilon$. So $\mathcal{U}R$ is an ultraring.\par
 
% END S2.SS2.p3
 
% BEGIN S2.SS2.p4
 
A ring homomorphism $h:R\to S$ induces a morphism of ultrarings $\mathcal{U}h:\mathcal{U}R\to\mathcal{U}S$ by applying $h$ pointwise to matrices. It is well known (see for instance [7]) that every symmetric monoidal, finite-coproduct-preserving functor $\mathcal{U}R\to\mathcal{U}S$ is monoidally isomorphic to a functor of the form $\mathcal{U}h$ for some $h:R\to S$. Therefore, if we consider the category of commutative rings $\mathbf{CRing}$ to be a $2$-category in which every hom-category is discrete, this defines a $2$-functor $\mathcal{U}:\mathbf{CRing}\to\mathbf{URing}$, and we have:\par
 
% END S2.SS2.p4
 
% BEGIN Thmtheorem4
 
\par\medskip\noindent\textbf{Proposition 4.}\par
 
% BEGIN Thmtheorem4.p1
 
$\mathcal{U}:\mathbf{CRing}\hookrightarrow\mathbf{URing}$ is fully faithful.\par
 
% END Thmtheorem4.p1
 
% END Thmtheorem4
 
% BEGIN S2.SS2.p5
 
Let $\mathbf{Fin}_{\ast}$ be the category of finite sets and partial functions. Observe that the cartesian product induces a symmetric monoidal structure on $\mathbf{Fin}_{\ast}$ (which is not a categorical product), whose unit is the singleton set $\{\ast\}$, and that $\mathbf{Fin}_{\ast}$ is obviously distributive (coproducts are disjoint unions) and has a zero object (the empty set). A Frobenius sesquimonoid structure may be defined by letting $\delta_{A}:A\to A\times A$ be the diagonal, $\varepsilon_{A}:A\to\{\ast\}$ be the unique total function and $\delta^{\ast}_{A}:A\times A\to A$ be the function which is defined on $(a,a^{\prime})$ exactly when $a=a^{\prime}$.\par
 
% END S2.SS2.p5
 
% BEGIN S2.SS2.p6
 
We invite the reader to check that, if $f:C\to A$ and $g:C\to B$ are partial functions, the condition $(f\otimes g)\delta=0$ amounts to the domains of $f$ and $g$ being disjoint. In that case, the pairing $\langle f,g\rangle:C\to A\oplus B$ is the evident partial function whose domain is the union of the domains of $f$ and $g$. Furthermore, observe that any partial function $p:A\to\{\ast\}$ satisfies $(p\otimes p)\delta=p$, and corresponds to a subset $P\subseteq A$. The partial function $\overline{p}:A\to\{\ast\}$ corresponding to the complement of $P$ is easily seen to be the unique one verifying the desired properties of point (2) of Definition 3.\par
 
% END S2.SS2.p6
 
% BEGIN S2.SS2.p7
 
The above shows that $\mathbf{Fin}_{\ast}$ has the structure of an ultraring, except that it is not small. This is easily fixed: let $\mathbb{F}_{1}$ be the category whose objects are natural numbers and whose arrows $n\to m$ are partial functions $\{1,\ldots,n\}\to\{1,\ldots,m\}$, or equivalently, $m\times n$ matrices with coefficients in $\{0,1\}$ such that $1$ appears at most once in each column. As a skeleton of $\mathbf{Fin}_{\ast}$, $\mathbb{F}_{1}$ is an ultraring. In fact, it is the initial ultraring: since $\mathbb{F}_{1}$ is generated under finite coproducts by the monoidal unit $1$ and its identity arrow, any finite-coproduct-preserving functor $F:\mathbb{F}_{1}\to\mathcal{R}$ with $\mathcal{R}$ an ultraring is determined by the image of $1$; since morphisms of ultrarings are monoidal, $F1\cong I$ and $\mathbf{URing}(\mathbb{F}_{1},\mathcal{R})$ is equivalent to the terminal category.\par
 
% END S2.SS2.p7
 
% BEGIN S2.SS2.p8
 
Observe that $\mathbf{Fin}_{\ast}$ is the Kleisli category of the partiality monad $-+\{\ast\}$ over the category $\mathbf{Fin}$ of finite sets and total functions. This is a special case of a more general situation. A category is lextensive if it has finite limits and disjoint, pullback-stable finite coproducts. A lextensive category is called Boolean if every subobject has a complement; see [8]. A primordial example is $\mathbf{Fin}$. In fact, for what concerns finite limits and finite coproducts, the objects of Boolean lextensive categories may be thought of as sets and their arrows as total functions. In particular, if $\mathcal{B}$ is a Boolean lextensive category whose terminal object is $1$, the endofunctor $-+1$ is a monad and the Kleisli category $\mathcal{B}_{\ast}$ may again be seen as a category of “sets and partial functions”. The Kleisli categories of the partiality monad on small Boolean lextensive categories may be characterized as certain ultrarings, which we introduce below.\par
 
% END S2.SS2.p8
 
% BEGIN Thmtheorem5
 
\par\medskip\noindent\textbf{Definition 5.}\par
 
% BEGIN Thmtheorem5.p1
 
An ultraring is Boolean if every arrow $f$ satisfies $\delta f=(f\otimes f)\delta$ and if for every $p:A\to I$ there exists a unique coproduct diagram $(\iota_{p},\overline{\iota}_{p})$ over $A$ such that $p=\varepsilon\iota_{p}^{\ast}$, $\overline{\iota}_{p}=\iota_{\overline{p}}$ and, for every arrow $f$, $\iota_{\varepsilon f}=(\mathrm{id}\otimes\varepsilon)\iota_{\varepsilon\delta^%
{\ast}(f\otimes\mathrm{id})}$.\par
 
% END Thmtheorem5.p1
 
% END Thmtheorem5
 
% BEGIN S2.SS2.p9
 
In the following, we denote by $\mathbf{BoolURing}$ the full sub-$2$-category of $\mathbf{URing}$ on Boolean ultrarings, and by $\mathbf{BoolLext}$ the $2$-category of small Boolean lextensive categories, finitely continuous functors which preserve finite coproducts, and natural transformations.\par
 
% END S2.SS2.p9
 
% BEGIN Thmtheorem6
 
\par\medskip\noindent\textbf{Proposition 6.}\par
 
% BEGIN Thmtheorem6.p1
 
The categories $\mathbf{BoolURing}$ and $\mathbf{BoolLext}$ are equivalent.\par
 
% END Thmtheorem6.p1
 
% END Thmtheorem6
 
% BEGIN S2.SS2.SSS0.P0.SPx1
 
\par\medskip\noindent\textbf{Proof.}\par
 
% BEGIN S2.SS2.SSS0.P0.SPx1.p1
 
We omit this proof due to space constraints. $\hfill\blacktriangleleft$\par
 
% END S2.SS2.SSS0.P0.SPx1.p1
 
% END S2.SS2.SSS0.P0.SPx1
 
% END S2.SS2
 
% END S2


% BEGIN S3
 
\par\medskip\noindent\textbf{3 Annular Theories and Models}\par
 
% BEGIN S3.SS1
 
\par\medskip\noindent\textbf{3.1 Presentations of Ultrarings}\par
 
% BEGIN S3.SS1.p1
 
Having introduced the motivating examples of ultrarings, we now wish to provide presentations of these objects, which will subsequently enable us to perform essential constructions on them, as well as providing the basis for a syntax having a natural semantics in ultrarings.\par
 
% END S3.SS1.p1
 
% BEGIN Thmtheorem7
 
\par\medskip\noindent\textbf{Definition 7.}\par
 
% BEGIN Thmtheorem7.p1
 
A signature $\mathfrak{S}$ is a pair $(\mathrm{Sort}(\mathfrak{S}),\mathrm{Gen}(\mathfrak{S}))$ such that:\par
  
\par\noindent  $\blacksquare$ 
% BEGIN S3.I1.i1.p1
 
$\mathrm{Sort}(\mathfrak{S})$ is a set of sorts. A tensor of arity $n$ is a sequence of $n$ sorts $A_{1}\otimes\cdots\otimes A_{n}$. The unique tensor of arity $0$ is denoted by $I$.\par
 
% END S3.I1.i1.p1
 
 
\par\noindent  $\blacksquare$ 
% BEGIN S3.I1.i2.p1
 
$\mathrm{Gen}(\mathfrak{S})$ is a set of generators, each with a type $g:T\to\bigoplus_{j=1}^{k}U_{j}$ where $T,U_{j}$ are tensors.\par
 
% END S3.I1.i2.p1
 
  
% END Thmtheorem7.p1
 
% END Thmtheorem7
 
% BEGIN S3.SS1.p2
 
We fix a countably infinite set of variables, ranged over by $x,y,z$… and we use $\vec{x},\vec{y}$,…to denote finite sequences of variables. Given a signature $\mathfrak{S}$, the monomials over $\mathfrak{S}$ are defined as follows:\par
    \[p,q::=(x=_{A}y)\mathrel{|}(g^{j}(\vec{x})=\vec{y})\mathrel{|}1\mathrel{|}%
\overline{p}\mathrel{|}pq\mathrel{|}{\smallint\!}_{x}p,\]    
where $A$ ranges over $\mathrm{Sort}(\mathfrak{S})$, $g:T\to\bigoplus_{j=1}^{k}U_{j}$ ranges over $\mathrm{Gen}(\mathfrak{S})$, $1\leq j\leq k$ and the length of $\vec{x}$ and $\vec{y}$ matches that of $T$ and $U_{j}$, respectively. We write $g^{j}(\vec{x})$, $(g^{j}=\vec{y})$ and $g^{j}$ rather than $(g^{j}(\vec{x})=)$, $(g^{j}()=\vec{y})$ and $(g^{j}()=)$, respectively. We abbreviate $(x_{1}=_{A_{1}}y_{1})\cdots(x_{n}=_{A_{n}}y_{n})$ to $(\vec{x}=_{\vec{A}}\vec{y})$. The notation ${\smallint\!}_{x}a$ is a binder: $x$ is bound in $a$ and is subject to the usual renaming conventions. If $\vec{x}=x_{1},\ldots,x_{n}$, we write ${\smallint\!}_{\vec{x}}a$ for ${\smallint\!}_{x_{1}}\cdots{\smallint\!}_{x_{n}}a$. When $\vec{x}$ is empty, ${\smallint\!}_{\vec{x}}a$ is just $a$. We call a monomial positive if it contains no instances of the negation operator $\overline{p}$.\par
 
% END S3.SS1.p2
 
% BEGIN S3.SS1.p3
 
Monomials generalize logical formulas. Intuitively, in the Boolean case, the monomials $1$, $\overline{p}$, $pq$ and ${\smallint\!}_{x}p$ could be written $\top$, $\lnot p$, $p\land q$, $\exists x.p$, and the monomial $(g^{j}(\vec{x})=\vec{y})$ corresponds to a relation symbol $g^{j}(\vec{x},\vec{y})$ satisfying that $g^{j}(\vec{x},\vec{y})\land g^{j}(\vec{x},\vec{z})$ implies $\vec{y}=\vec{z}$, i.e., a functional relation. The equality monomial $(x=_{A}y)$ is already a standard Boolean formula.\par
 
% END S3.SS1.p3
 
% BEGIN S3.SS1.p4
 
A pre-polynomial is a formal finite sum of monomials, with $0$ representing the empty sum. We can extend the operations of multiplication and integration in the definition of monomials (the last two operations) to pre-polynomials by assuming distributivity and linearity, respectively.\par
 
% END S3.SS1.p4
 
% BEGIN S3.SS1.p5
 
A context is a finite list of type assignments of the form $x:A$ where $x$ is a variable and $A$ a sort, such that no variable appears twice. If $\vec{x}=x_{1},\ldots,x_{n}$ is repetition-free and $T=\bigotimes_{i=1}^{n}A_{i}$, we also write the context $x_{1}:A_{1},\ldots,x_{n}:A_{n}$ as $\vec{x}:T$. In particular, in $\vec{x}:I$ the sequence $\vec{x}$ is empty. If $\Gamma$ and $\Delta$ are contexts, $\Gamma,\Delta$ denotes their concatenation (the variables appearing in each are assumed disjoint).\par
 
% END S3.SS1.p5
 
% BEGIN S3.F1
  
\begin{center}\includegraphics[width=0.95\linewidth,height=0.45\textheight,keepaspectratio]{graphic-01.pdf}\end{center}
  
\begin{center}\includegraphics[width=0.95\linewidth,height=0.45\textheight,keepaspectratio]{graphic-02.pdf}\end{center}
  
Figure 1: Monomials in context. We omit a rule allowing to permute type assignments in contexts.\par
 
% END S3.F1
 
% BEGIN S3.SS1.p6
 
A monomial-in-context is an expression of the form $\Gamma\vdash p::\Delta$, where $p$ is a monomial and $\Gamma$ and $\Delta$ are contexts, which is derivable by means of the inductive rules of Fig. 1. If $\Gamma\vdash p::\Delta$ is derivable, we say that $p$ is well-typed, that $\Gamma$ is a domain for $p$ and that the context $\Gamma,\Delta$ matches $p$. Note that, thanks to the penultimate rule of Fig. 1, any context matching $p$ is also a domain for it. Pre-polynomials-in-context are defined similarly, adding a rule to derive $\Gamma\vdash\sum_{i\in I}p_{i}::\Delta$ whenever $\Gamma\vdash p_{i}::\Delta$ is derivable for all $i\in I$ (in particular, $\Gamma\vdash 0::\Delta$ is derivable for any $\Gamma,\Delta$). From now on, we will only consider well-typed monomials and pre-polynomials.\par
 
% END S3.SS1.p6
 
% BEGIN S3.F2
    $\displaystyle{\smallint\!}_{x}{\smallint\!}_{y}p$ $\displaystyle\approx_{\Gamma}{\smallint\!}_{y}{\smallint\!}_{x}p$ $\displaystyle{\smallint\!}_{y}(x=y)$ $\displaystyle\approx_{\Gamma,x:A}1$     $\displaystyle(x=_{A}y)p$ $\displaystyle\approx_{\Gamma,x:A,y:A}(x=_{A}y)p[x/y]$ $\displaystyle(x=_{A}x)$ $\displaystyle\approx_{\Gamma,x:A}1$     $\displaystyle(x=_{A}y)$ $\displaystyle\approx_{\Gamma,x:A,y:A}(y=_{A}x)$ $\displaystyle p+\overline{p}$ $\displaystyle\approx_{\Gamma}1$     $\displaystyle{\smallint\!}_{x}pq$ $\displaystyle\approx_{\Gamma}p{\smallint\!}_{x}q\quad\text{if $x$ is not free %
in $p$}$    
Figure 2: Basic equations for congruences on pre-polynomials. $p[x/y]$ denotes usual substitution.\par
 
% END S3.F2
 
% BEGIN S3.SS1.p7
 
A congruence is a context-indexed family of equivalence relations $\approx\ :=(\approx_{\Gamma})_{\Gamma}$ on pre-polynomials matching $\Gamma$, closed under associativity and commutativity of multiplication $pq$, neutrality of $1$, every equation of Fig. 2 and such that:\par
  
\par\noindent  $\blacksquare$ 
% BEGIN S3.I2.i1.p1
 
$p\approx_{\Gamma}q$ implies $\overline{p}\approx_{\Gamma}\overline{q}$, $pr\approx_{\Gamma}qr$ for all $r$ matching $\Gamma$ and $p\approx_{\Gamma^{\prime}}q$ for every context $\Gamma^{\prime}$ including the assignments of $\Gamma$, possibly in a different order;\par
 
% END S3.I2.i1.p1
 
 
\par\noindent  $\blacksquare$ 
% BEGIN S3.I2.i2.p1
 
$p\approx_{\Gamma,x:A}q$ implies ${\smallint\!}_{x}p\approx_{\Gamma}{\smallint\!}_{x}q$;\par
 
% END S3.I2.i2.p1
 
 
\par\noindent  $\blacksquare$ 
% BEGIN S3.I2.i3.p1
 
$p\approx_{\Gamma}q$ iff $p+r\approx_{\Gamma}q+r$ for any $r$ matching $\Gamma$ (we call this cancellativity).\par
 
% END S3.I2.i3.p1
 
  
Given a set $S$ of equations of the form $p\approx_{\Gamma}q$ where $p,q$ are pre-polynomials matching $\Gamma$, the congruence generated by $S$ is the smallest congruence containing the equations of $S$.\par
 
% END S3.SS1.p7
 
% BEGIN S3.SS1.p8
 
In the sequel, if $\approx$ is a congruence and we write $p\approx q$, we mean $p\approx_{\Gamma}q$ for every $\Gamma$ matching both $p,q$.\par
 
% END S3.SS1.p8
 
% BEGIN Thmtheorem8
 
\par\medskip\noindent\textbf{Lemma 8.}\par
 
% BEGIN Thmtheorem8.p1
 
For any congruence $\approx$ and any pre-polynomial $p$, $p^{2}\approx p$ iff $p\overline{p}\approx 0$. Moreover, either condition implies $\overline{p}^{2}\approx\overline{p}$.\par
 
% END Thmtheorem8.p1
 
% END Thmtheorem8
 
% BEGIN S3.SS1.SSS0.P0.SPx1
 
\par\medskip\noindent\textbf{Proof.}\par
 
% BEGIN S3.SS1.SSS0.P0.SPx1.p1
 
See §A.1. $\hfill\blacktriangleleft$\par
 
% END S3.SS1.SSS0.P0.SPx1.p1
 
% BEGIN S3.F3
  
\begin{center}\includegraphics[width=0.95\linewidth,height=0.45\textheight,keepaspectratio]{graphic-03.pdf}\end{center}
  
\begin{center}\includegraphics[width=0.95\linewidth,height=0.45\textheight,keepaspectratio]{graphic-04.pdf}\end{center}
  
Figure 3: Compatibility of pre-polynomials. The side condition $(*)$ is: $\Gamma,\vec{y}:U,\vec{z}:V$ (for some $U,V$) are domains for $r,s,t$, respectively, and $\vec{x}$ are the outputs of $r$.\par
 
% END S3.F3
 
% BEGIN S3.SS1.SSS0.P0.SPx1.p2
 
Given a congruence $\approx$ over a signature $\mathfrak{S}$, we define $\approx$-compatibility to be the context-indexed family of smallest symmetric relations $\frown_{\Gamma}$ on monomials of $\mathfrak{S}$ closed under the rules of Fig. 3 such that $p\frown_{\Gamma}q$ implies that $\Gamma$ is a domain for both $p$ and $q$. We write $p\frown q$ to mean $p\frown_{\Gamma}q$ for every $\Gamma$ which is a domain for both $p$ and $q$.\par
 
% END S3.SS1.SSS0.P0.SPx1.p2
 
% BEGIN S3.SS1.SSS0.P0.SPx1.p3
 
A $\approx$-polynomial (or simply polynomial when this is unambiguous) is a pre-polynomial $p$ such that there exists $\Gamma$ such that $q\frown_{\Gamma}r$ for all monomials $q,r$ of $p$ and such that, whenever the expression $\overline{q}$ appears in $p$, we have $q^{2}\approx q$.\par
 
% END S3.SS1.SSS0.P0.SPx1.p3
 
% BEGIN Thmtheorem9
 
\par\medskip\noindent\textbf{Definition 9.}\par
 
% BEGIN Thmtheorem9.p1
 
An annular theory $\mathfrak{T}$ (which we abbreviate to theory in the remainder) is a triple $(\mathrm{Sort}(\mathfrak{T}),\mathrm{Gen}(\mathfrak{T}),\mathrm{Eq}(\mathfrak{%
T}))$ where $(\mathrm{Sort}(\mathfrak{T}),\mathrm{Gen}(\mathfrak{T}))$ is a signature and $\mathrm{Eq}(\mathfrak{T})$ is a well-founded set of equations of the form $p\approx_{\Gamma}q$ where $p,q$ are pre-polynomials matching $\Gamma$, such that, calling such an equation $e$, we have that $p$ and $q$ are $\approx$-polynomials with $\approx$ the congruence generated by the equations of $\mathrm{Eq}(\mathfrak{T})$ strictly below $e$ in the order. We denote by $\approx^{\mathfrak{T}}$ the congruence generated by $\mathrm{Eq}(\mathfrak{T})$.\par
 
% END Thmtheorem9.p1
 
% END Thmtheorem9
 
% BEGIN Thmtheorem10
 
\par\medskip\noindent\textbf{Definition 10.}\par
 
% BEGIN Thmtheorem10.p1
 
A $\approx$-matrix (or simply matrix) of type $\bigoplus_{k\in K}\vec{x}^{k}:T_{k}\to\bigoplus_{j\in J}\vec{y}^{j}:U_{j}$ consists of a $(J\times K)$-indexed family of $\approx$-polynomials-in-context $\left(\vec{x}^{k}:T_{k}\vdash p_{jk}::\vec{y}^{j}:U_{j}\right)_{(j,k)\in J%
\times K}$ which we abbreviate to $(p_{jk})$, such that,\par
    \[\text{for every $j\neq j^{\prime}\in J$, $k\in K$, and monomials $r$ in $p_{jk%
}$ and $r^{\prime}$ in $p_{j^{\prime}k}$, $r\frown_{\vec{x}^{k}:T_{k}}r^{%
\prime}$.}\hfill({\dagger})\]    
We always identify $\approx$-matrices whose entries are equal modulo $\approx$.\par
 
% END Thmtheorem10.p1
 
% END Thmtheorem10
 
% BEGIN Thmtheorem11
 
\par\medskip\noindent\textbf{Lemma 11.}\par
 
% BEGIN Thmtheorem11.p1
 
Let $(p_{jk})$ and $(q_{hj})$ be $\approx$-matrices of type $\bigoplus_{k\in K}\vec{x}^{k}:T_{k}\to\bigoplus_{j\in J}\vec{y}^{j}:U_{j}$ and $\bigoplus_{j\in J}\vec{y}^{j}:U_{j}\to\bigoplus_{h\in H}\vec{z}^{h}:V_{h}$, respectively. Then, the family of pre-polynomials-in-context indexed by $H\times J$ defined by\par
    \[(q\circ p)_{hk}:=\left(\vec{x}^{k}:T_{k}\vdash\sum_{j\in J}{\smallint\!}_{\vec%
{y}^{j}}q_{hj}p_{jk}::\vec{z}^{h}:V_{h}\right)_{(h,k)\in H\times K}\]    
is a $\approx$-matrix of type $\bigoplus_{k\in K}\vec{x}^{k}:T_{k}\to\bigoplus_{h\in H}\vec{z}^{h}:V_{h}$.\par
 
% END Thmtheorem11.p1
 
% END Thmtheorem11
 
% END S3.SS1.SSS0.P0.SPx1
 
% BEGIN S3.SS1.SSS0.P0.SPx2
 
\par\medskip\noindent\textbf{Proof.}\par
 
% BEGIN S3.SS1.SSS0.P0.SPx2.p1
 
See §A.1. $\hfill\blacktriangleleft$\par
 
% END S3.SS1.SSS0.P0.SPx2.p1
 
% BEGIN S3.SS1.SSS0.P0.SPx2.p2
 
Let $\mathfrak{T}$ be a theory. We define a category $\mathbb{F}_{1}[\mathfrak{T}]$ as follows:\par
  
\par\noindent  $\blacksquare$ 
% BEGIN S3.I3.i1.p1
 
the objects are formal finite coproducts $\bigoplus_{j\in J}\vec{x}^{j}:T_{j}$ of contexts of $\mathfrak{T}$, modulo injective renaming of variables.\par
 
% END S3.I3.i1.p1
 
 
\par\noindent  $\blacksquare$ 
% BEGIN S3.I3.i2.p1
 
If $A:=\bigoplus_{k\in K}\vec{x}^{k}:T_{k}$ and $B:=\bigoplus_{j\in J}\vec{y}^{j}:U_{j}$, $\mathbb{F}_{1}[\mathfrak{T}](A,B)$ is the set of $\approx^{\mathfrak{T}}$-matrices of type $A\to B$.\par
 
% END S3.I3.i2.p1
 
 
\par\noindent  $\blacksquare$ 
% BEGIN S3.I3.i3.p1
 
Composition is given by Lemma 11, and $\mathrm{id}_{A}$ is the diagonal matrix indexed by $K\times K$ which, at position $(k,k)$, is equal to $(\vec{x}^{k}=_{T_{k}}\vec{z}^{k})$, where we used the fact that, since objects are defined up to variable renaming, by renaming each $x^{k}_{h}$ to $z^{k}_{h}$, we may write $A$ as $\bigoplus_{k\in K}\vec{z}^{k}:T_{k}$.\par
 
% END S3.I3.i3.p1
 
  
As defined, $\mathbb{F}_{1}[\mathfrak{T}]$ is not small, because the finite sets indexing the objects are arbitrary. This may be circumvented by restricting to finite subsets of a fixed, countably infinite set of indices. We leave this implicit for notational convenience, but treat $\mathbb{F}_{1}[\mathfrak{T}]$ as a small category.\par
 
% END S3.SS1.SSS0.P0.SPx2.p2
 
% BEGIN Thmtheorem12
 
\par\medskip\noindent\textbf{Proposition 12.}\par
 
% BEGIN Thmtheorem12.p1
 
$\mathbb{F}_{1}[\mathfrak{T}]$ is an ultraring.\par
 
% END Thmtheorem12.p1
 
% END Thmtheorem12
 
% END S3.SS1.SSS0.P0.SPx2
 
% BEGIN S3.SS1.SSS0.P0.SPx3
 
\par\medskip\noindent\textbf{Proof.}\par
 
% BEGIN S3.SS1.SSS0.P0.SPx3.p1
 
We spell out the required structures and properties, without entering into the details. Here “polynomial”, “matrix”, etc. mean $\approx^{\mathfrak{T}}$-polynomial, $\approx^{\mathfrak{T}}$-matrix, etc.\par
 
% END S3.SS1.SSS0.P0.SPx3.p1
 
% BEGIN S3.SS1.SSS0.P0.SPx3.p2
 
Let $A:=\bigoplus_{j\in J}\vec{x}^{j}.T_{j}$ and $B:=\bigoplus_{k\in K}\vec{y}^{k}.U_{k}$ be generic objects of $\mathbb{F}_{1}[\mathfrak{T}]$. The (strict) symmetric monoidal structure of $\mathbb{F}_{1}[\mathfrak{T}]$ is given by $A\otimes B:=\bigoplus_{(j,k)\in J\times K}\vec{x}^{j}\vec{y}^{k}.T_{j}\otimes U%
_{k}$ on objects, with $I$ as unit. On morphisms, tensor is matrix product, valid by a similar argument to the proof of Lemma 11, but without the integrals. The empty sum $0$ is a zero object, the all-zero matrix of suitable size is the zero morphism between any two objects. The coproduct of $A$ and $B$ (where we may suppose $J\cap K=\emptyset$ and $\vec{x}^{j},\vec{y}^{k}$ to be disjoint) is given by $\bigoplus_{l\in J\cup K}\vec{z}^{l}.V_{l}$, where $\vec{z}^{l}$ and $V_{l}$ are equal to $\vec{x}^{j}$ and $T_{j}$ (resp. $\vec{y}^{k}$ and $U_{k}$) if $l\in J$ (resp. $l\in K$). This obviously satisfies strict distributivity. The injections $\iota_{1},\iota_{2}$ are the block column matrices $\left(\begin{array}[]{c}\mathrm{id}\\
0\end{array}\right)$ and $\left(\begin{array}[]{c}0\\
\mathrm{id}\end{array}\right)$, respectively, whereas the copairing of two matrices $P,Q$ is the block row matrix $(P\ Q)$.\par
 
% END S3.SS1.SSS0.P0.SPx3.p2
 
% BEGIN S3.SS1.SSS0.P0.SPx3.p3
 
For what concerns the Frobenius sesquimonoid structure, $\delta_{A}$ is the matrix indexed over $J^{2}\times J$ whose column for $j\in J$ has the element $(\vec{x}^{j}=\vec{y}^{j})(\vec{x}^{j}=\vec{z}^{j})$ at position $(j,j)$ and is zero everywhere else, where we are renaming variables so that $A\otimes A=\bigoplus_{j,j^{\prime}\in J}\vec{y}^{j}\vec{z}^{j^{\prime}}.T_{j}%
\otimes T_{j^{\prime}}$. $\varepsilon_{A}$ is the row matrix which is equal to the monomial $1$ everywhere. $\delta^{\ast}_{A}$ is the transpose of $\delta_{A}$.\par
 
% END S3.SS1.SSS0.P0.SPx3.p3
 
% BEGIN S3.SS1.SSS0.P0.SPx3.p4
 
Thanks to the above definition, elementary calculations show that, if $P=(p_{kj})_{j\in J,k\in K}$ and $Q=(q_{lj})_{j\in J,l\in L}$ are matrices corresponding to morphisms $A\to B$ and $A\to C$ for some other object $C$ (indexed by $L$, which we may supposed to be disjoint from $K$ even if $C=B$), then $(P\otimes Q)\delta_{A}\approx^{\mathfrak{T}}0$ is equivalent to asking that, for all $j\in J$, $k\in K$ and $l\in L$, we have $p_{kj}q_{lj}\approx^{\mathfrak{T}}0$, which implies $p_{kj}\frown q_{lj}$, so we may form the block column matrix $R:=\left(\begin{array}[]{c}P\\
Q\end{array}\right)$ (with $K\cup L$ rows and $J$ columns). The reader may check that the coinjections are the block row matrices $(\mathrm{id}\ 0)$ and $(0\ \mathrm{id})$, so $R$ is the (unique) pairing of $P,Q$, as desired.\par
 
% END S3.SS1.SSS0.P0.SPx3.p4
 
% BEGIN S3.SS1.SSS0.P0.SPx3.p5
 
Let now $P=(p_{k})_{k\in K}$ be a matrix of type $A\to I$. This means that $P$ is a row matrix, and elementary calculations show that $(P\otimes P)\delta\approx^{\mathfrak{T}}P$ is equivalent to asking that $p_{k}^{2}\approx^{\mathfrak{T}}p_{k}$ for all $k\in K$, so we are allowed to consider the polynomials $\overline{p_{k}}$ and define $\overline{P}:=(\overline{p_{k}})_{k\in K}$. We need to show that $(P\otimes\overline{P})\delta\approx^{\mathfrak{T}}0$. By similar calculations as above, this is equivalent to $p_{k}\overline{p_{k}}\approx^{\mathfrak{T}}0$ for all $k\in K$, which holds by Lemma 8. We are left with proving that $\nabla\langle P,\overline{P}\rangle\approx^{\mathfrak{T}}\varepsilon$. Notice that, in general, $\nabla\langle Q,R\rangle=Q+R$ as matrices. This sum exists because, as shown above, $\langle Q,R\rangle$ is the block column matrix $\left(\begin{array}[]{c}Q\\
R\end{array}\right)$. So we need to show that $p_{k}+\overline{p_{k}}\approx^{\mathfrak{T}}1$ for all $k\in K$, but this is one of the equations of Fig. 2. $\hfill\blacktriangleleft$\par
 
% END S3.SS1.SSS0.P0.SPx3.p5
 
% END S3.SS1.SSS0.P0.SPx3
 
% END S3.SS1
 
% BEGIN S3.SS2
 
\par\medskip\noindent\textbf{3.2 Models and the Canonical Theory}\par
 
% BEGIN S3.SS2.p1
 
The term “annular theory” in Definition 9 is not merely an analogy with logic: it is chosen so that ultraring morphisms out of $\mathbb{F}_{1}[\mathfrak{T}]$ become models of $\mathfrak{T}$, in a sense that we now establish.\par
 
% END S3.SS2.p1
 
% BEGIN Thmtheorem13
 
\par\medskip\noindent\textbf{Definition 13.}\par
 
% BEGIN Thmtheorem13.p1
 
Given a signature $\mathfrak{S}$ and an ultraring $\mathcal{A}$, an $\mathfrak{S}$-structure in $\mathcal{A}$ consists of:\par
  
\par\noindent  $\blacksquare$ 
% BEGIN S3.I4.i1.p1
 
An assignment $M_{0}:\mathrm{Sort}(\mathfrak{S})\to\mathrm{ob}(\mathcal{A})$. This extends to an assignment on tensors given by $M_{0}(\bigotimes_{i}A_{i}):=\bigotimes_{i}M_{0}(A_{i})$.\par
 
% END S3.I4.i1.p1
 
 
\par\noindent  $\blacksquare$ 
% BEGIN S3.I4.i2.p1
 
An assignment $M_{1}:\mathrm{Gen}(\mathfrak{S})\to\mathrm{mor}(\mathcal{A})$ such that if $g:T\to U$ then $M_{1}(g):M_{0}(T)\to M_{0}(U)$.\par
 
% END S3.I4.i2.p1
 
  
For $\mathfrak{S}$-structures $M,M^{\prime}$ in $\mathcal{A}$, a homomorphism $m:M\to M^{\prime}$ consists of a collection of morphisms $m_{X}:M_{0}(X)\to M^{\prime}_{0}(X)$ in $\mathcal{A}$ indexed by $X\in\mathrm{Sort}(\mathfrak{S})$ such that for each $g:\bigotimes_{i}A_{i}\to\bigotimes_{j}B_{j}\in\mathrm{Gen}(\mathfrak{S})$, $M_{1}(g);\bigotimes_{j}m_{B_{j}}=\bigotimes_{i}m_{A_{i}};M^{\prime}_{1}(g)$.\par
 
% END Thmtheorem13.p1
 
% END Thmtheorem13
 
% BEGIN S3.SS2.p2
 
Given a theory $\mathfrak{T}$ over a signature $\mathfrak{S}$ and a $\mathfrak{S}$-structure $M$, we can define an interpretation over $M$ of any $\approx^{\mathfrak{T}}$-polynomial-in-context, as detailed in §A.2. The notion of model is then defined in the expected way:\par
 
% END S3.SS2.p2
 
% BEGIN Thmtheorem14
 
\par\medskip\noindent\textbf{Definition 14.}\par
 
% BEGIN Thmtheorem14.p1
 
Let $\mathfrak{T}$ be a theory over signature $\mathfrak{S}$ and $\mathcal{A}$ an ultraring. A model of $\mathfrak{T}$ in $\mathcal{A}$ is a $\mathfrak{S}$-structure $M$ in $\mathcal{A}$ such that for each equation $p\approx_{\Gamma,\Delta}q$ of $\mathfrak{T}$, we have $M(\Gamma\vdash p::\Delta)=M(\Gamma\vdash q::\Delta)$.\par
 
% END Thmtheorem14.p1
 
% BEGIN Thmtheorem14.p2
 
For models $M,M^{\prime}$ of $\mathfrak{T}$ in $\mathcal{A}$, a homomorphism $M\to M^{\prime}$ is simply a homomorphism of $\mathfrak{S}$-structures. We denote by $\mathfrak{T}\text{-}\mathrm{Mod}(\mathcal{A})$ the category of models of $\mathfrak{T}$ in $\mathcal{A}$ and homomorphisms between them.\par
 
% END Thmtheorem14.p2
 
% END Thmtheorem14
 
% BEGIN Thmtheorem15
 
\par\medskip\noindent\textbf{Proposition 15.}\par
 
% BEGIN Thmtheorem15.p1
 
Let $\mathfrak{T}$ be a theory. There is an $\mathcal{A}$-natural equivalence of categories $\mathfrak{T}\text{-}\mathrm{Mod}(\mathcal{A})\simeq\mathbf{URing}(\mathbb{F}_{%
1}[\mathfrak{T}],\mathcal{A}).$\par
 
% END Thmtheorem15.p1
 
% END Thmtheorem15
 
% BEGIN S3.SS2.SSS0.P0.SPx1
 
\par\medskip\noindent\textbf{Proof.}\par
 
% BEGIN S3.SS2.SSS0.P0.SPx1.p1
 
See §A.2. $\hfill\blacktriangleleft$\par
 
% END S3.SS2.SSS0.P0.SPx1.p1
 
% BEGIN S3.SS2.SSS0.P0.SPx1.p2
 
Proposition 15 implies in particular the existence of a generic model of any annular theory $\mathfrak{T}$, corresponding to the identity functor on $\mathbb{F}_{1}[\mathfrak{T}]$, of which a model in any ultraring is an image by naturality.\par
 
% END S3.SS2.SSS0.P0.SPx1.p2
 
% BEGIN S3.SS2.SSS0.P0.SPx1.p3
 
If Proposition 12 shows us that we can generate ultrarings from theories, it may not be a surprise that every ultraring arises in this way up to equivalence.\par
 
% END S3.SS2.SSS0.P0.SPx1.p3
 
% BEGIN Thmtheorem16
 
\par\medskip\noindent\textbf{Definition 16.}\par
 
% BEGIN Thmtheorem16.p1
 
We define the canonical theory of an ultraring $\mathcal{R}$ to be the theory $\mathfrak{T}$ where:\par
  
\par\noindent  $\blacksquare$ 
% BEGIN S3.I5.i1.p1
 
$\mathrm{Sort}(\mathfrak{T})=\mathrm{ob}(\mathcal{R})$, where we denote the sort corresponding to $A$ by $\ulcorner A\urcorner$;\par
 
% END S3.I5.i1.p1
 
 
\par\noindent  $\blacksquare$ 
% BEGIN S3.I5.i2.p1
 
$\mathrm{Gen}(\mathfrak{T})$ has a generator $\ulcorner g\urcorner:\bigotimes_{j\in J}\ulcorner A_{j}\urcorner\to\bigoplus_{%
k\in K}\bigotimes_{l\in L_{k}}\ulcorner B_{k,l}\urcorner$ for each morphism $g:\bigotimes_{j\in J}A_{j}\to\bigoplus_{k\in K}\bigotimes_{l\in L_{k}}B_{k,l}$ in $\mathcal{R}$,333This is a mild abuse of notation, since the name $\ulcorner g\urcorner$ can refer to multiple generators depending on how the domain and codomain are decomposed. When not specified, the type will be the one with the minimal decomposition.\par
 
% END S3.I5.i2.p1
 
 
\par\noindent  $\blacksquare$ 
% BEGIN S3.I5.i3.p1
 
$\mathrm{Eq}(\mathfrak{T})$ consists of equations:\par
  
\par\noindent  – 
% BEGIN S3.I5.i3.I1.i1.p1
 
$\left(\ulcorner\mathrm{id}\urcorner(\vec{x})=\vec{y}\right)\approx(\vec{x}=%
\vec{y})$ for $\ulcorner\mathrm{id}\urcorner:\bigotimes_{j\in J}\ulcorner A_{i}\urcorner\to%
\bigotimes_{j\in J}\ulcorner A_{i}\urcorner$,\par
 
% END S3.I5.i3.I1.i1.p1
 
 
\par\noindent  – 
% BEGIN S3.I5.i3.I1.i2.p1
 
$\int_{\vec{y}}(\ulcorner f\urcorner(\vec{x})=\vec{y})(\ulcorner g\urcorner(%
\vec{y})=\vec{z})\approx(\ulcorner gf\urcorner(\vec{x})=\vec{z})$ for composable $f,g$ in $\mathcal{R}$ (with no coproduct decomposition of the codomains).\par
 
% END S3.I5.i3.I1.i2.p1
 
 
\par\noindent  – 
% BEGIN S3.I5.i3.I1.i3.p1
 
$(\ulcorner\iota^{*}_{j}g\urcorner(\vec{x})=\vec{y})\approx(\ulcorner g%
\urcorner^{j}(\vec{x})=\vec{y})$.\par
 
% END S3.I5.i3.I1.i3.p1
 
 
\par\noindent  – 
% BEGIN S3.I5.i3.I1.i4.p1
 
$(\ulcorner f\otimes g\urcorner(\vec{x},\vec{u})=\vec{y},\vec{v})\approx(%
\ulcorner f\urcorner(\vec{x})=\vec{y})(\ulcorner g\urcorner(\vec{u})=\vec{v})$.\par
 
% END S3.I5.i3.I1.i4.p1
 
 
\par\noindent  – 
% BEGIN S3.I5.i3.I1.i5.p1
 
$(\ulcorner 0\urcorner(\vec{x})=\vec{y})\approx_{\vec{x}:T,\vec{y}:U,\Gamma}0$, for any typing of $\ulcorner 0\urcorner$.\par
 
% END S3.I5.i3.I1.i5.p1
 
 
\par\noindent  – 
% BEGIN S3.I5.i3.I1.i6.p1
 
$(\ulcorner\delta\urcorner(\vec{x})=\vec{x_{1}},\vec{x_{2}})\approx(\vec{x}=%
\vec{x}_{1})(\vec{x}=\vec{x}_{2})$\par
 
% END S3.I5.i3.I1.i6.p1
 
 
\par\noindent  – 
% BEGIN S3.I5.i3.I1.i7.p1
 
$\ulcorner\varepsilon\urcorner(\vec{x})\approx_{\vec{x}:T,\Gamma}1$\par
 
% END S3.I5.i3.I1.i7.p1
 
 
\par\noindent  – 
% BEGIN S3.I5.i3.I1.i8.p1
 
$(\ulcorner\delta^{*}\urcorner(\vec{x_{1}},\vec{x_{2}})=\vec{x})\approx(\vec{x}%
_{1}=\vec{x})(\vec{x}_{2}=\vec{x})$.\par
 
% END S3.I5.i3.I1.i8.p1
 
  
% END S3.I5.i3.p1
 
  
% END Thmtheorem16.p1
 
% END Thmtheorem16
 
% BEGIN Thmtheorem17
 
\par\medskip\noindent\textbf{Theorem 17.}\par
 
% BEGIN Thmtheorem17.p1
 
Let $\mathcal{R}$ be an ultraring and $\mathfrak{T}$ its canonical theory, as defined in Definition 16. Then $\mathcal{R}\simeq\mathbb{F}_{1}[\mathfrak{T}]$.\par
 
% END Thmtheorem17.p1
 
% END Thmtheorem17
 
% END S3.SS2.SSS0.P0.SPx1
 
% BEGIN S3.SS2.SSS0.P0.SPx2
 
\par\medskip\noindent\textbf{Proof.}\par
 
% BEGIN S3.SS2.SSS0.P0.SPx2.p1
 
See §A.2. $\hfill\blacktriangleleft$\par
 
% END S3.SS2.SSS0.P0.SPx2.p1
 
% END S3.SS2.SSS0.P0.SPx2
 
% END S3.SS2
 
% BEGIN S3.SS3
 
\par\medskip\noindent\textbf{3.3 Examples of Annular Theories}\par
 
% BEGIN S3.SS3.p1
 
In what follows, we let $\mathbb{K}$ denote either $\mathbb{F}_{1}$ or the ultraring corresponding to a commutative ring $R$ with no non-trivial idempotent, i.e., in which $r^{2}=r$ implies $r\in\{0,1\}$ for all $r\in R$. Examples include any field or, more generally, any integral domain such as $\mathbb{Z}$.\par
 
% END S3.SS3.p1
 
% BEGIN S3.SS3.p2
 
We define a predicate of a theory $\mathfrak{T}$ to be a $\approx^{\mathfrak{T}}$-polynomial $p:T\to I$ such that $p^{2}\approx_{\vec{x}:T}^{\mathfrak{T}}p$. Observe that, if $M$ is a model of $\mathfrak{T}$ in $\mathbb{K}$, then $M(p)$ must be a row vector of zeros and ones: the equation $p^{2}\approx_{\vec{x}:T}^{\mathfrak{T}}p$ becomes $(M(p)\otimes M(p))\delta_{T}=M(p)$, and we already observed that this is equivalent to all entries of $M(p)$ being idempotent. Furthermore, $M(T)=I^{\oplus n}$ for some $n$ (all objects of $\mathbb{K}$ are of this form). If we see $M(T)$ as the set $\{1,\ldots,n\}$, $M(p)$ may be identified with the subset of $M(T)$ of all $i$ such $M(p)_{i}=1$, justifying the terminology.\par
 
% END S3.SS3.p2
 
% BEGIN S3.SS3.p3
 
Predicates are closed under Boolean logic: $0$ and $1$ are predicates; if $p$ and $q$ are predicates, then $pq$ and $\overline{p}$ are predicates (Lemma 8). Also, by the second row of Fig. 2, $(x=_{A}y)$ is always a predicate: $(x=_{A}y)^{2}\approx_{x:A,y:A}(x=_{A}y)(x=_{A}x)\approx_{x:A,y:A}(x=y)$. Moreover, models in $\mathbb{K}$ are consistent with Boolean logic: if $M$ and $p$ are as above, $M(0)=\emptyset$, $M(1)=M(T)$, $M(pq)=M(p)\cap M(q)$, $M(\overline{p})=M(T)\setminus M(p)$ and $M(\vec{x}=_{T}\vec{y})$ is the diagonal of $M(T)\times M(T)$.\par
 
% END S3.SS3.p3
 
% BEGIN S3.SS3.p4
 
In the sequel, when we say that a generator $P$ is a predicate over $T$, we mean that $P:T\to I$ and that the equation $P^{2}\approx_{\vec{x}:T}^{\mathfrak{T}}P$ is added to the theory. By the above, a monomial $\varphi$ consisting entirely of generators which are predicates is like a Boolean formula and, if $\vec{x}:T$ matches $\varphi$, we may impose that $\forall\vec{x}.\varphi$ holds by adding the equation $\varphi\approx_{\vec{x}:T}1$.\par
 
% END S3.SS3.p4
 
% BEGIN S3.SS3.p5
 
For example, we define the theory $\mathfrak{Ord}$ with one sort $N$ and one predicate $\leq$ over $N\otimes N$ with equations making $\leq$ is a total order. We write $(x\leq y)$ rather than $\mathop{\leq}(x,y)$. A model of $\mathfrak{Ord}$ in $\mathbb{K}$ is precisely a finite total order; modulo iso, we may assume it to be of the form $[n]:=\{1<\cdots<n\}$. The theory $\mathfrak{Str}$ is obtained from $\mathfrak{Ord}$ by adding a generator $X:N\to I$ with no further equations. A model $s$ of $\mathfrak{Str}$ in $\mathbb{K}$ is a finite total order indexing a row vector $s(X)$ of zeros and ones if $\mathbb{K}=\mathbb{F}_{1}$, or of elements of $R$ otherwise. In other words, $s$ represents a binary string when $\mathbb{K}=\mathbb{F}_{1}$, or a string of elements of $R$ otherwise. We denote the $i$-th bit/element by $s(X)_{i}$.\par
 
% END S3.SS3.p5
 
% END S3.SS3
 
% END S3


% BEGIN A1.SS1
 
\par\medskip\noindent\textbf{A.1 Proofs of Some Lemmas}\par
 
% BEGIN A1.SS1.SSS0.P0.SPx1
 
\par\medskip\noindent\textbf{Proof of Lemma 2.}\par
 
% BEGIN A1.SS1.SSS0.P0.SPx1.p1
 
We start by observing that, if $\iota_{i}:B_{i}\to B_{1}\oplus B_{2}$ are injections, then $\iota_{i}\delta^{\ast}_{B_{i}}(\iota^{\ast}_{i}\otimes\mathrm{id}_{B_{i}})=%
\delta^{\ast}_{B_{1}\oplus B_{2}}(\mathrm{id}_{B_{1}\oplus B_{2}}\otimes\iota_%
{i})$. We prove this for $i=1$, the other case is essentially identical. We omit obvious subscripts for identity arrows. Let $\chi=[\iota_{1}\otimes\mathrm{id},\iota_{2}\otimes\mathrm{id}]:(B_{1}\otimes B%
_{1})\oplus(B_{2}\otimes B_{1})\to(B_{1}\oplus B_{2})\otimes B_{1}$ be the distributivity isomorphism. We have, using (1), $\iota_{1}\delta^{\ast}_{B_{1}}(\iota^{\ast}_{1}\otimes\mathrm{id})\chi=\iota_{%
1}\delta^{\ast}_{B_{1}}[\iota^{\ast}_{1}\iota_{1}\otimes\mathrm{id},\iota^{%
\ast}_{1}\iota_{2}\otimes\mathrm{id}]=\iota_{1}\delta^{\ast}_{B_{1}}[\mathrm{%
id},0]=[\iota_{1}\delta^{\ast}_{B_{1}},0]=[\delta^{\ast}_{B_{1}\oplus B_{2}}(%
\iota_{1}\otimes\iota_{1}),\delta^{\ast}_{B_{1}\oplus B_{2}}(\iota_{2}\otimes%
\iota_{1})]=\delta^{\ast}_{B_{1}\oplus B_{2}}(\mathrm{id}\otimes\iota_{1})\chi$, which proves the claim because $\chi$ is invertible. We will also need that, if $\kappa_{i}:A\otimes B_{i}\to(A\otimes B_{1})\oplus(A\otimes B_{2})$ are injections and $\chi^{\prime}:(A\otimes B_{1})\oplus(A\otimes B_{2})\to A\otimes(B_{1}\oplus B%
_{2})$ another instance of distributivity, then $\chi^{\prime}\kappa_{i}=\mathrm{id}\otimes\iota_{i}$, which holds by definition of $\chi^{\prime}$.\par
 
% END A1.SS1.SSS0.P0.SPx1.p1
 
% BEGIN A1.SS1.SSS0.P0.SPx1.p2
 
Suppose now that we have $f,g:A\rightrightarrows B_{1}\oplus B_{2}$ with $\iota_{i}^{*}f=\iota_{i}^{*}g$ for $i=1,2$. Using the above observations, we have $\delta^{\ast}_{B_{1}\oplus B_{2}}(f\otimes\mathrm{id})\chi^{\prime}\kappa_{i}=%
\delta^{\ast}_{B_{1}\oplus B_{2}}(f\otimes\iota_{i})=\iota_{i}\delta^{\ast}_{B%
_{i}}(\iota^{\ast}_{i}f\otimes\mathrm{id})=\iota_{i}\delta^{\ast}_{B_{i}}(%
\iota^{\ast}_{i}g\otimes\mathrm{id})$, from which, applying the first two equalities in reverse order, we obtain $\delta^{\ast}_{B_{1}\oplus B_{2}}(f\otimes\mathrm{id})\chi^{\prime}\kappa_{i}=%
\delta^{\ast}_{B_{1}\oplus B_{2}}(g\otimes\mathrm{id})\chi^{\prime}\kappa_{i}$ for $i=1,2$. Since injections are jointly epic and $\chi^{\prime}$ is invertible, we infer $\delta^{\ast}_{B_{1}\oplus B_{2}}(f\otimes\mathrm{id})=\delta^{\ast}_{B_{1}%
\oplus B_{2}}(g\otimes\mathrm{id})$, from which we conclude $f=g$ by definition $\delta^{\ast}$-category. $\hfill\blacktriangleleft$\par
 
% END A1.SS1.SSS0.P0.SPx1.p2
 
% END A1.SS1.SSS0.P0.SPx1
 
% BEGIN A1.SS1.SSS0.P0.SPx2
 
\par\medskip\noindent\textbf{Proof of Lemma 8.}\par
 
% BEGIN A1.SS1.SSS0.P0.SPx2.p1
 
We start by observing that cancellativity implies $\overline{\overline{p}}\approx p$, as well as $\overline{1}\approx 0$. Next, we claim that, whenever $q$ matches the same contexts as $p$ and $q^{2}\approx q$, we have $\overline{pq}\approx\overline{p}q+p\overline{q}+\overline{p}\,\overline{q}$. Indeed, $pq+\overline{pq}\approx 1\approx(p+\overline{p})(q+\overline{q})\approx pq+%
\overline{p}q+p\overline{q}+\overline{p}\,\overline{q}$, so we conclude by cancellativity.\par
 
% END A1.SS1.SSS0.P0.SPx2.p1
 
% BEGIN A1.SS1.SSS0.P0.SPx2.p2
 
Suppose that $p\overline{p}\approx 0$. Then, $p\approx p(p+\overline{p})\approx p^{2}+p\overline{p}\approx p^{2}$. For the converse, if $p^{2}\approx p$, then using the above claim we have $\overline{p\overline{p}}=p^{2}+\overline{p}^{2}+\overline{p}p=p+\overline{p}(%
\overline{p}+p)=p+\overline{p}=1$, so $p\overline{p}=0$ by the first two observations. Finally, we have $\overline{p}\approx\overline{p^{2}}\approx p\overline{p}+\overline{p}p+%
\overline{p}^{2}\approx\overline{p}^{2}$. $\hfill\blacktriangleleft$\par
 
% END A1.SS1.SSS0.P0.SPx2.p2
 
% END A1.SS1.SSS0.P0.SPx2
 
% BEGIN A1.SS1.SSS0.P0.SPx3
 
\par\medskip\noindent\textbf{Proof of Lemma 11.}\par
 
% BEGIN A1.SS1.SSS0.P0.SPx3.p1
 
For indices $h,j,j^{\prime},k$, let $r,r^{\prime},s,s^{\prime}$ be monomials in $p_{j}k,p_{j^{\prime}k},q_{hj},q_{hj^{\prime}}$, respectively. We need to check that each $(q\circ p)_{hk}$ is a polynomial by verifying that, when $j\neq j^{\prime}$ we have ${\smallint\!}_{\vec{y}^{j}}sr\frown_{\vec{x}^{k}:T_{k}}{\smallint\!}_{\vec{y}^%
{j^{\prime}}}s^{\prime}r^{\prime}$. By hypothesis, $r\frown_{\vec{x}^{k}:T_{k}}r^{\prime}$, so we conclude by the last rule of Fig. 3. We must also check that the context is valid. By hypothesis, $q_{hj}p_{jk}$ is matched by $\vec{x}^{k}:T_{k},\vec{y}^{j}:U_{j},\vec{z}^{h}:V_{h}$, with $\vec{y}^{j}$ and $\vec{z}^{h}$ being outputs, hence each monomial ${\smallint\!}_{\vec{y}^{j}}rs$ is well-typed and matched by $\vec{x}^{k}:T_{k},\vec{z}^{h}:V_{h}$, as required. Finally, we must check condition (${\dagger}$) of Definition 10. Given $h\neq h^{\prime}$, we require ${\smallint\!}_{\vec{y}^{j}}sr\frown_{\vec{x}^{k}:T_{k}}{\smallint\!}_{\vec{y}^%
{j^{\prime}}}t^{\prime}r^{\prime}$, where $t^{\prime}$ is a monomial of $q_{h^{\prime}j^{\prime}}$. If $j\neq j^{\prime}$, we conclude as above. If $j=j^{\prime}$, by hypothesis we have $s\frown_{\vec{y}^{j}:U_{j}}t^{\prime}$, so we conclude by applying the last rule of Fig. 3 to these instead. $\hfill\blacktriangleleft$\par
 
% END A1.SS1.SSS0.P0.SPx3.p1
 
% END A1.SS1.SSS0.P0.SPx3
 
% END A1.SS1


% BEGIN A1.SS2
 
\par\medskip\noindent\textbf{A.2 Interpreting Polynomials in a $\mathfrak{S}$-structure}\par
 
% BEGIN A1.SS2.p1
 
Let $M$ be a $\mathfrak{S}$-structure. We start by inductively defining an interpretation over $M$ for positive monomials-in-context, following Fig. 1.\par
  
\par\noindent  $\blacksquare$ 
% BEGIN A1.I1.i1.p1
 
As a base case, we interpret monomials over $\mathfrak{S}$ with the maximal codomain context of all outputs and the minimal domain of all remaining free variables, for each of the basic constructions of monomials:\par
  
\par\noindent  – 
% BEGIN A1.I1.i1.I1.i1.p1
 
$M\left(x:A\vdash(x=_{A}y)::y:A\right)$ is $\mathrm{id}_{M_{0}(A)}$,\par
 
% END A1.I1.i1.I1.i1.p1
 
 
\par\noindent  – 
% BEGIN A1.I1.i1.I1.i2.p1
 
$M\left(\vec{x}:T\vdash(g^{j}(\vec{x})=\vec{y})::\vec{y}:U\right)$ is $M_{1}(g);\iota^{*}_{j}$,\par
 
% END A1.I1.i1.I1.i2.p1
 
 
\par\noindent  – 
% BEGIN A1.I1.i1.I1.i3.p1
 
$M\left(\vdash 1::\right)$ is $\varepsilon_{I}=\mathrm{id}_{I}$,\par
 
% END A1.I1.i1.I1.i3.p1
 
  
% END A1.I1.i1.p1
 
 
\par\noindent  $\blacksquare$ 
% BEGIN A1.I1.i2.p1
 
Next, we describe the effect of expanding and reorganizing contexts.\par
  
\par\noindent  – 
% BEGIN A1.I1.i2.I1.i1.p1
 
Given $f:=M(\Gamma\vdash p::\Delta,x:A)$, let $M(\Gamma,x:A\vdash p::\Delta):=\varepsilon_{M_{0}(A)}\delta^{*}_{M_{0}(A)}(f%
\otimes\mathrm{id}_{A})$.\par
 
% END A1.I1.i2.I1.i1.p1
 
 
\par\noindent  – 
% BEGIN A1.I1.i2.I1.i2.p1
 
Relatedly, given $f:=M(\Gamma\vdash p::\Delta)$, we define $M(x:A,\Gamma\vdash p::\Delta)$ to be $\beta(\varepsilon_{M_{0}(A)}\otimes f)$, (recalling that $\beta:I\otimes I\to I$ is the canonical isomorphism).\par
 
% END A1.I1.i2.I1.i2.p1
 
 
\par\noindent  – 
% BEGIN A1.I1.i2.I1.i3.p1
 
Given $M(\Gamma\vdash p::\Delta)$, permutations of $\Gamma$ and $\Delta$ are interpreted by pre- or post-composing with the corresponding symmetry morphisms.\par
 
% END A1.I1.i2.I1.i3.p1
 
  
% END A1.I1.i2.p1
 
 
\par\noindent  $\blacksquare$ 
% BEGIN A1.I1.i3.p1
 
Finally, we specify how monomials combine.\par
  
\par\noindent  – 
% BEGIN A1.I1.i3.I1.i1.p1
 
Given $f:=M(\Gamma\vdash p::\Delta,\vec{x}:T)$ and $g:=M(\vec{x}:T,\Gamma^{\prime}\vdash q::\Delta^{\prime})$ we can define $M\left(\Gamma,\Gamma^{\prime}\vdash\int_{\vec{x}:T}pq::\Delta,\Delta^{\prime}\right)$ to be $(\mathrm{id}\otimes g)(f\otimes\mathrm{id})$; to recover an interpretation of $\int_{\vec{x}}p$, we identify this with $\int_{\vec{x}}p1$.\par
 
% END A1.I1.i3.I1.i1.p1
 
 
\par\noindent  – 
% BEGIN A1.I1.i3.I1.i2.p1
 
Given $f:=M(x_{1}:A,x_{2}:A,\Gamma\vdash p::\Delta)$ and a variable $x$ not appearing in $p$, we define $M(x:A,\Gamma\vdash p[x_{1}/x][x_{2}/x]::\Delta)$ to be $f(\delta_{M_{0}(A)}\otimes\mathrm{id})$.\par
 
% END A1.I1.i3.I1.i2.p1
 
 
\par\noindent  – 
% BEGIN A1.I1.i3.I1.i3.p1
 
Dually, given $f:=M(\Gamma\vdash p::x_{1}:B,x_{2}:B,\Delta)$ and a variable $x$ not appearing in $p$, we define $M(\Gamma\vdash p[y_{1}/y][y_{2}/y]::y:B,\Delta)$ to be $(\delta^{*}_{M_{0}(B)}\otimes\mathrm{id})f$.\par
 
% END A1.I1.i3.I1.i3.p1
 
  
% END A1.I1.i3.p1
 
  
% END A1.SS2.p1
 
% BEGIN A1.SS2.p2
 
Let us check that the above interpretations are well-defined. The first set of rules give unambiguous interpretations. For the second set, there is exactly one way to arrive a given context using these rules up to reordering, so it suffices to check that the order of operations does not affect the result; this is straightforward, relying on the fact that $\delta^{*}$ is symmetric.\par
 
% END A1.SS2.p2
 
% BEGIN A1.SS2.p3
 
The only place where ambiguity is introduced in the third set of rules, where the interpretation ostensibly depends on how the domain and codomain contexts are partitioned. We take these case by case.\par
  
\par\noindent  $\blacksquare$ 
% BEGIN A1.I2.i1.p1
 
Consider $M\left(\Gamma,\Gamma^{\prime}\vdash\int_{\vec{x}:T}pq::\Delta,\Delta^{\prime}\right)$. By disjointness, $\Delta$ and $\Delta^{\prime}$ must respectively contain only outputs of $p$ and $q$, so this is unambiguous; the remaining free variables of $p$ and $q$ must respectively go into $\Gamma$ and $\Gamma^{\prime}$. Any remaining free variables, whether assigned to $\Gamma$ or $\Gamma^{\prime}$, must be interpreted by $\varepsilon$, which eliminates any remaining ambiguity.\par
 
% END A1.I2.i1.p1
 
 
\par\noindent  $\blacksquare$ 
% BEGIN A1.I2.i2.p1
 
In $M(x:A,\Gamma\vdash p[x_{1}/x][x_{2}/x]::\Delta)$, there is formally a choice to make about how to assign the instances of $x$ to $x_{1}$ or $x_{2}$; however, since the inductive construction will eventually disambiguate between any pair of repeated instances, commutativity and associativity of $\delta$ guarantee that any such choice will result in the same interpretation (considering also that the interpretation is invariant under $\alpha$-equivalence).\par
 
% END A1.I2.i2.p1
 
 
\par\noindent  $\blacksquare$ 
% BEGIN A1.I2.i3.p1
 
Similarly, distinguishing duplicates of variables in the codomain produces a well-defined result by commutativity and associativity of $\delta^{*}$, so we are done.\par
 
% END A1.I2.i3.p1
 
  
% END A1.SS2.p3
 
% BEGIN A1.SS2.p4
 
Next, given $f:=M\left(\Gamma\vdash p::\right)$, we have that $M\left(\Gamma\vdash p^{2}::\right)$ is exactly $(f\otimes f)\delta$. As such, when these interpretations are equal, we may take $M\left(\Gamma\vdash\overline{p}::\right)$ to be $\overline{f}$ provided by Definition 3, thus extending the above to monomials featuring $\overline{p}$.\par
 
% END A1.SS2.p4
 
% BEGIN A1.SS2.p5
 
Finally, we define $M\left(\Gamma\vdash 0::\Delta\right)$ to be the zero morphism of the appropriate type.\par
 
% END A1.SS2.p5
 
% BEGIN Thmtheorem23
 
\par\medskip\noindent\textbf{Lemma 23.}\par
 
% BEGIN Thmtheorem23.p1
 
Suppose we are given a theory $\mathfrak{T}$ over a signature $\mathfrak{S}$ and a $\mathfrak{S}$-structure $M$ in $\mathcal{A}$. Suppose that whenever we have positive monomials-in-context $\Gamma\vdash p::\Delta$ and $\Gamma\vdash q::\Delta$ such that $p\approx_{\Gamma,\Delta}q$, it is the case that $M(\Gamma\vdash p::\Delta)=M(\Gamma\vdash q::\Delta)$, so that in particular $M$ can interpret all monomials which are $\approx$-polynomials. Suppose further that $M$ identifies $p$ and $q$ whenever $p\approx q$ when $p$ and $q$ are either monomials or $0$. Then $p\frown_{\mathfrak{T}}q$ implies the pairing $\langle p,q\rangle$ exists in $\mathcal{A}$.\par
 
% END Thmtheorem23.p1
 
% END Thmtheorem23
 
% BEGIN A1.SS2.SSS0.P0.SPx1
 
\par\medskip\noindent\textbf{Proof.}\par
 
% BEGIN A1.SS2.SSS0.P0.SPx1.p1
 
It suffices to check the axioms of Fig. 3 which determine the compatibility relation. Indeed, $pq\approx 0$ guarantees that if $M$ respectively interprets $p,q$ as $f,g$ in a given context then the first axiom of ultrarings Definition 3 ensures that the pairing exists.\par
 
% END A1.SS2.SSS0.P0.SPx1.p1
 
% BEGIN A1.SS2.SSS0.P0.SPx1.p2
 
The pairing of two components $j,j^{\prime}$ of (the interpretation of) a generator $g$ is the result of coinjecting $g$ onto the coproduct of their codomains.\par
 
% END A1.SS2.SSS0.P0.SPx1.p2
 
% BEGIN A1.SS2.SSS0.P0.SPx1.p3
 
If $p\frown q$ and $\langle p,q\rangle$ exists, we can obtain the pairing of the interpretations of ${\smallint\!}_{\vec{x},\vec{y}}spr$ and ${\smallint\!}_{\vec{x},\vec{z}}tqr$ by precomposing with $r$ and postcomposing with $s\oplus t$.\par
 
% END A1.SS2.SSS0.P0.SPx1.p3
 
% BEGIN A1.SS2.SSS0.P0.SPx1.p4
 
The existence of pairings in the last instance is immediate, so we are done. $\hfill\blacktriangleleft$\par
 
% END A1.SS2.SSS0.P0.SPx1.p4
 
% BEGIN A1.SS2.SSS0.P0.SPx1.p5
 
As such, given $\mathfrak{T}$ and a model $M$ satisfying the conditions of Lemma 23, $M$ has a well-defined interpretation of any $\approx$-polynomial-in-context $\Gamma\vdash p::\Delta$ as the joint pairing of the monomials appearing in $p$ (with the same context), followed by the codiagonal map.\par
 
% END A1.SS2.SSS0.P0.SPx1.p5
 
% BEGIN A1.SS2.SSS0.P0.SPx1.p6
 
% BEGIN Thmtheorem15
 
\par\medskip\noindent\textbf{Proposition 15. [Restated, see original statement.]}\par
 
% BEGIN Thmtheorem15.p1
 
Let $\mathfrak{T}$ be a theory. There is an $\mathcal{A}$-natural equivalence of categories $\mathfrak{T}\text{-}\mathrm{Mod}(\mathcal{A})\simeq\mathbf{URing}(\mathbb{F}_{%
1}[\mathfrak{T}],\mathcal{A}).$\par
 
% END Thmtheorem15.p1
 
% END Thmtheorem15
 
% END A1.SS2.SSS0.P0.SPx1.p6
 
% END A1.SS2.SSS0.P0.SPx1
 
% BEGIN A1.SS2.SSS0.P0.SPx2
 
\par\medskip\noindent\textbf{Proof.}\par
 
% BEGIN A1.SS2.SSS0.P0.SPx2.p1
 
We have thus far shown that a $\mathfrak{T}$-model $M$ in $\mathcal{A}$ uniquely determines interpretations of $\approx^{\mathfrak{T}}$-polynomials-in-context. To produce a functor $\mathbb{F}_{1}[\mathfrak{T}]\to\mathcal{A}$, we need to extend this to $\approx^{\mathfrak{T}}$-matrices. Each entry of a matrix $(p_{jk})$ determines a morphism $M_{0}(T_{k})\to M_{0}(U_{j})$. Condition (${\dagger}$) of Definition 10 ensures that for fixed $k$, the interpretations of $p_{jk}$ can be paired to give a morphism $M_{0}(T_{k})\to\bigoplus_{j\in J}M_{0}(U_{j})$; via the universal property of the coproduct, these produce a morphism $\bigoplus_{k\in K}M_{0}(T_{k})\to\bigoplus_{j\in J}M_{0}(U_{j})$. This extended construction ensures that the monoidal product and coproduct, as well as $\delta$, $\delta^{*}$ and $\epsilon$, are preserved.\par
 
% END A1.SS2.SSS0.P0.SPx2.p1
 
% BEGIN A1.SS2.SSS0.P0.SPx2.p2
 
Conversely, a morphism of ultrarings $F:\mathbb{F}_{1}[\mathfrak{T}]\to\mathcal{A}$ defines a model $M$ by its restriction to the sorts and generators. The equations of $\mathfrak{T}$ necessarily hold because they correspond to morphisms forced to be equal in $\mathbb{F}_{1}[\mathfrak{T}]$, so are mapped by $F$ to the equal morphisms in $\mathcal{A}$.\par
 
% END A1.SS2.SSS0.P0.SPx2.p2
 
% BEGIN A1.SS2.SSS0.P0.SPx2.p3
 
It is straightforward to verify that a model morphism uniquely determines a monoidal natural transformation, and conversely, since the latter are determined by their components at objects which generate the category under monoidal products and coproducts.\par
 
% END A1.SS2.SSS0.P0.SPx2.p3
 
% BEGIN A1.SS2.SSS0.P0.SPx2.p4
 
Finally, naturality in $\mathcal{A}$ follows from the fact that morphisms of ultrafunctors preserve the required structure: composing with an ultraring morphism $\mathcal{A}\to\mathcal{B}$ turns a $\mathfrak{T}$-model in $\mathcal{A}$ into one in $\mathcal{B}$. $\hfill\blacktriangleleft$\par
 
% END A1.SS2.SSS0.P0.SPx2.p4
 
% BEGIN A1.SS2.SSS0.P0.SPx2.p5
 
% BEGIN Thmtheorem17
 
\par\medskip\noindent\textbf{Theorem 17. [Restated, see original statement.]}\par
 
% BEGIN Thmtheorem17.p1
 
Let $\mathcal{R}$ be an ultraring and $\mathfrak{T}$ its canonical theory, as defined in Definition 16. Then $\mathcal{R}\simeq\mathbb{F}_{1}[\mathfrak{T}]$.\par
 
% END Thmtheorem17.p1
 
% END Thmtheorem17
 
% END A1.SS2.SSS0.P0.SPx2.p5
 
% END A1.SS2.SSS0.P0.SPx2
 
% BEGIN A1.SS2.SSS0.P0.SPx3
 
\par\medskip\noindent\textbf{Proof.}\par
 
% BEGIN A1.SS2.SSS0.P0.SPx3.p1
 
We have a functor $\ulcorner-\urcorner:\mathcal{R}\to\mathbb{F}_{1}[\mathfrak{T}]$ sending each object $X$ to the sort $\ulcorner X\urcorner$ and each morphism $g:X\to Y$ to the corresponding generator (with no decomposition of $X$ and $Y$ as tensors or coproducts). We show that $\ulcorner-\urcorner$ is an equivalence of ultrarings.\par
 
% END A1.SS2.SSS0.P0.SPx3.p1
 
% BEGIN A1.SS2.SSS0.P0.SPx3.p2
 
Consider the generators (which we distinguish by a subscript) $\ulcorner\mathrm{id}\urcorner_{1}:\ulcorner\bigotimes_{j\in J}A_{j}\urcorner%
\to\bigotimes_{j\in J}\ulcorner A_{j}\urcorner$ and $\ulcorner\mathrm{id}\urcorner_{2}:\bigotimes_{j\in J}\ulcorner A_{j}\urcorner%
\to\ulcorner\bigotimes_{j\in J}A_{j}\urcorner$. By the second and first equations, these are mutual inverses in $\mathbb{F}_{1}[\mathfrak{T}]$. Moreover, naturality of these morphisms follows from the fourth axiom (applied inductively), so $\ulcorner-\urcorner$ preserves the monoidal product.\par
 
% END A1.SS2.SSS0.P0.SPx3.p2
 
% BEGIN A1.SS2.SSS0.P0.SPx3.p3
 
Consider the generator $\ulcorner\mathrm{id}\urcorner:\ulcorner\bigoplus_{j\in J}A_{j}\urcorner\to%
\bigoplus_{j\in J}\ulcorner A_{j}\urcorner$. By the third equation, we have that $\left(\ulcorner\mathrm{id}\urcorner^{j}(\vec{x})=\vec{y}\right)\approx(%
\ulcorner\iota^{*}_{j}\urcorner(\vec{x})=\vec{y}),$ so this is equal to the row vector $(\ulcorner\iota^{*}_{j_{1}}\urcorner,\dotsc,\ulcorner\iota^{*}_{j_{n}}\urcorner)$.\par
 
% END A1.SS2.SSS0.P0.SPx3.p3
 
% BEGIN A1.SS2.SSS0.P0.SPx3.p4
 
Dually, we have generators $\ulcorner\iota_{i}\urcorner:\ulcorner A_{i}\urcorner\to\bigoplus_{j\in J}%
\ulcorner A_{j}\urcorner$. By the third, first and fifth equations, we have $\left(\ulcorner\iota_{i}\urcorner^{j}(\vec{x})=\vec{y}\right)\approx\left(%
\ulcorner\iota^{*}_{j}\iota_{i}\urcorner(\vec{x})=\vec{y}\right)$ which gives $(\vec{x}=\vec{y})$ if $i=j$ and $0$ otherwise. Thus $\ulcorner-\urcorner$ preserves the coproduct injections; in particular, these morphisms are compatible, so we can construct the column vector $(\ulcorner\iota_{j_{1}}\urcorner,\dotsc,\ulcorner\iota_{j_{n}}\urcorner)^{t}$. We can use the second and third axioms to conclude that this column vector is the inverse to $\ulcorner\mathrm{id}\urcorner$, so $\ulcorner-\urcorner$ preserves coproducts. Formally, we must also check that the universal morphisms are respected by this isomorphism, in the sense that $\int_{\vec{y}}\left(\ulcorner[f_{1},\dotsc,f_{n}]\urcorner(\vec{x})=\vec{y}%
\right)(\ulcorner\mathrm{id}\urcorner(\vec{y})=\vec{z})\approx\int_{\vec{y}^{%
\prime}}(\ulcorner\mathrm{id}\urcorner(\vec{x})=\vec{y}^{\prime})\left(\left(%
\ulcorner f_{1}\urcorner,\dotsc,\ulcorner f_{n}\urcorner\right)(\vec{y}^{%
\prime})=\vec{z}\right).$ It suffices to check this componentwise by composing with $\iota_{i}$, which works by the preceding observations and the second axiom.\par
 
% END A1.SS2.SSS0.P0.SPx3.p4
 
% BEGIN A1.SS2.SSS0.P0.SPx3.p5
 
For the $\delta^{*}$-category structure, we observe that the sixth, seventh and eighth axioms directly ensure that these are preserved up to the isomorphisms constructed above. With this, we have shown that $\ulcorner-\urcorner$ is a valid morphism of ultrarings.\par
 
% END A1.SS2.SSS0.P0.SPx3.p5
 
% BEGIN A1.SS2.SSS0.P0.SPx3.p6
 
Conversely, there is a canonical model of the canonical theory $\mathfrak{T}$ in $\mathcal{R}$, given by interpreting $\ulcorner A\urcorner$ as $A$ and $\ulcorner g\urcorner$ as $g$ (with the relevant typing). The claim that this $\mathfrak{S}$-structure is indeed a model amounts to saying that the interpretations of the monomials appearing in the axioms of Definition 16 are equal in $\mathcal{R}$, which is straightforward to check.\par
 
% END A1.SS2.SSS0.P0.SPx3.p6
 
% BEGIN A1.SS2.SSS0.P0.SPx3.p7
 
Finally, to see that $\ulcorner-\urcorner$ is an equivalence, we simply observe on the one hand that the induced model in $\mathbb{F}_{1}[\mathfrak{T}]$ is isomorphic to that corresponding to the identity functor and on the other that the induced endofunctor of $\mathcal{R}$ is naturally isomorphic to the identity. The involved natural isomorphisms are the structural ones constructed above. $\hfill\blacktriangleleft$\par
 
% END A1.SS2.SSS0.P0.SPx3.p7
 
% END A1.SS2.SSS0.P0.SPx3
 
% END A1.SS2
